% =====================================================================
%  Shared CV body. Do not compile this file directly:
%  compile cv-ds.tex (data science) or cv-research.tex (research/PhD).
% =====================================================================
\providecommand{\cvversion}{ds}

\documentclass[9pt]{extarticle}
\usepackage[T1]{fontenc}
\usepackage[margin=0.3in]{geometry}
\usepackage{multicol}
\usepackage{titlesec}
\usepackage{csquotes}
\usepackage[dvipsnames]{xcolor}
\usepackage{tcolorbox}
\usepackage{enumitem}
\usepackage[symbol]{footmisc}
\usepackage{tikz}
\usetikzlibrary{calc}
\usepackage{tabularx}
\usepackage{etoolbox}
\usepackage[hidelinks, final]{hyperref}

\newtoggle{research}
\ifdefstring{\cvversion}{research}{\toggletrue{research}}{\togglefalse{research}}
\newcommand{\dsonly}[1]{\iftoggle{research}{}{#1}}
\newcommand{\rsonly}[1]{\iftoggle{research}{#1}{}}
\newcommand{\pick}[2]{\iftoggle{research}{#2}{#1}}

\definecolor{Headings}{HTML}{577359}
\definecolor{grey}{gray}{0.95}
\definecolor{hg}{HTML}{577359}

\titleformat{\section}{\Large\bfseries}{}{0em}{}[\color{Headings}{\titlerule[1.2pt]}]
\titleformat{\subsection}{\bfseries}{}{0em}{}

\setlist[itemize]{label={\color{gray}\scriptsize\textbullet}, nosep, itemsep=3pt, topsep=3pt, leftmargin=1em}

\renewcommand{\thefootnote}{\fnsymbol{footnote}}
\setlength{\parindent}{0pt}


\newcounter{cvcount}
\newcommand{\resetcvcount}{\setcounter{cvcount}{0}}

\newcommand{\timelinedot}{%
    \stepcounter{cvcount}%
    \begin{tikzpicture}[remember picture, overlay, baseline=-0.7ex]
      \node[circle, fill=hg, inner sep=1pt] (currentdot) {};
        \ifnum\value{cvcount}=1
            \node[coordinate] (tls) at (currentdot) {};
        \fi
        \node[coordinate] (tle) at (currentdot) {};
    \end{tikzpicture}%
}

\newcommand{\drawtimeline}{%
    \begin{tikzpicture}[remember picture, overlay]
        \draw[hg, line width=0.5pt] (tls) -- (tle);
    \end{tikzpicture}%
}

\newcommand{\cventry}[2]{
    \timelinedot &
    \textbf{#1} \\
    & \vspace{-4pt}
      \begin{itemize} #2 \end{itemize} \\[8pt]
}

\renewcommand{\maketitle}{
	\begin{center}
		{\huge\bfseries HENRY GROVES\\}
		\vspace{.5em}
		{\large hengro44@gmail.com | 07392 072855 | \href{https://github.com/hengrv}{github.com/hengrv} | \href{https://www.linkedin.com/in/hengrv/}{linkedin.com/in/hengrv}}
	\end{center}
}

\newcommand{\me}{\textbf{Groves, H.}}

\newcommand{\publist}{%
	\begin{itemize}
		\item \me, Jackson, L.F., Robertson, B.-A., and Hance, J.R.
		      \textit{Modelling Emotional Memory in Children with Tensor Networks}.
		      Preprint, \href{https://arxiv.org/abs/2606.28470}{arXiv:2606.28470}, 2026.
		\rsonly{\item Robertson, B.-A.*, Jackson, L., \me, and Hance, J.R.
		      \textit{Emotional Modulation of Temporal Order Memory in Childhood: A Quantum-Informed Account}.
		      Oral presentation, QIP26, Växjö, Sweden, June 2026. {\footnotesize *presenter}}
	\end{itemize}%
}

\begin{document}
\thispagestyle{empty}
\maketitle

\vspace{-.2cm}
\begin{minipage}{0.97\textwidth}
	\section{Education}
	\subsection{Newcastle University \hfill 2023--2027}
	\vspace{-.2cm}
	\begin{itemize}
		\item \textbf{M.Sc.\ + B.Sc.\ (Hons) Computer Science}
		\item \textbf{B.Sc.\ dissertation:} \textit{Quantum-Inspired Models of Emotional Memory}, mark 81 (First)%
		      \rsonly{. Supervised by Dr.~J.R.~Hance; developed into a co-authored preprint}
	\end{itemize}
	\vspace{-.4cm}
	\subsection{Pontefract New College \hfill 2021--2023}
	\vspace{-0.2cm}
	\begin{itemize}
		\item \textbf{A Levels:} Mathematics (\textbf{A*}), Further Mathematics (\textbf{A}), Computer Science (\textbf{A}), Physics (\textbf{A})
		\item \textbf{EPQ:} ``To what extent will quantum computing replace high performance computing in computational biology?'' (\textbf{A})
	\end{itemize}
	\dsonly{%
	\vspace{-.4cm}
	\subsection{Carleton High School, Pontefract \hfill 2016--2021}
	\vspace{-0.6cm}
	\begin{multicols}{3}
		\begin{itemize}
			\item \textbf{GCSEs:} 5 $\times$ grade 9, 3 $\times$ grade 8
			\item \textbf{BTECs:} \textbf{D*2}, \textbf{D2}
		\end{itemize}
	\end{multicols}
	}

\end{minipage}

\rsonly{%
	\section{Publications \& Presentations}
	\publist
	\section{Research Experience}
	\subsection{B.Sc.\ Dissertation, Quantum Group, Newcastle University \hfill 2025--2026}
	\vspace{-.1cm}
	\begin{itemize}
		\item Built quantum-inspired matrix product state (MPS) models of children's temporal order memory, encoding emotional valence and recall outcome as a joint composite-site index.
		\item Compared binary (qubit), three-state (qutrit) and valence-primed models with classical baselines (Naive Bayes, Markov chain) under 5-fold cross-validation; the valence-primed model reached 77.98\% accuracy.
		\item Implemented in Julia with ITensors.jl; extended into a co-authored preprint and a QIP26 conference talk.
	\end{itemize}
}

\vspace{.2cm}

\dsonly{\begin{minipage}[t]{0.65\textwidth}\vspace{0pt}}
	\raggedright

	\section{\pick{Work Experience}{Industry Experience}}
	\subsection{\large Northern Powergrid --- \normalsize\textit{Returning Summer Internship}}

	\resetcvcount
	\begin{tabularx}{\textwidth}{@{} >{\centering}p{1.5em} @{} X @{}}

		\cventry{\textit{Data Scientist (DSO)}\hfill June--September 2026}{
			\item \pick
			      {Designed a shared development framework for a growing data science team, built on Reproducible Analytical Pipeline (RAP) principles, with standardised logging, configuration and CI/CD on Azure DevOps.}
			      {Designed a shared development framework for the data science team, built on Reproducible Analytical Pipeline (RAP) principles, with standardised logging, configuration management and CI/CD.}
			\item Migrated 2 legacy manual Power BI/Excel processes to RAP pipelines, cutting run time from several days to a few minutes.
		}

		\cventry{\textit{Data Scientist (DSO)}\hfill May--September 2025}{
			\item \pick
			      {Built a graph-based simulation tool modelling high-voltage networks under normal and fault conditions, used across teams for network planning and forecasting low-carbon technology (LCT) uptake.}
			      {Developed a bottom-up, graph-based simulation of high-voltage networks under normal and fault conditions from historical and real-time asset data, replacing a top-down approach; used across teams for network planning.}
			\dsonly{\item Replaced a top-down approach with bottom-up simulation from historical and real-time asset data, enabling finer-grained analysis.}
			\item Modelled grid evolution to 2050 under 4 future energy scenarios.
		}

		\cventry{\textit{Data Scientist (DSO)}\hfill May--September 2024}{
			\item \pick
			      {Designed an end-to-end ML pipeline forecasting short-term load across $\sim$800 distribution substations from half-hourly meter data, with gap-filling via Akima interpolation and lag-based methods; used for network planning and procuring flexibility.}
			      {Built an end-to-end pipeline forecasting short-term load across $\sim$800 substations from half-hourly meter data, with Akima and lag-based gap-filling.}
			\item \pick
			      {Clustered load profiles with DTW k-means and showed that one model per cluster generalised across its members, reducing the number of models from $\sim$800 to 6.}
			      {Showed that substations clustered by DTW k-means share forecastable structure: one model per cluster generalised across all members, reducing $\sim$800 models to 6.}
			\item \pick
			      {Designed and tuned TCN and BiTCN neural networks for each cluster, with selected exogenous features improving accuracy.}
			      {Designed and tuned TCN and BiTCN models with exogenous features; results presented at CIRED 2025 as first author.}
		}

		\cventry{\textit{Data Analyst (System Forecasting)}\hfill May--August 2023}{
			\item Built a statistical demand model for $\sim$4M customers on low-voltage networks using SQL and Python, informing short- and long-term business planning.
			\dsonly{\item Prepared data for regulatory submissions under tight deadlines.}
		}
	\end{tabularx}

	\drawtimeline

	\vspace{-.7cm}

	\dsonly{%
	\section{Projects}
	\subsection{``Biggmarket'' --- item-swapping app promoting sustainability}
	Full-stack web app built with Next.js, TypeScript, tRPC and Tailwind CSS on PostgreSQL/Prisma, featuring a swipe interface, real-time messaging, Google OAuth and a user review system.
	}

\dsonly{\end{minipage}\hspace{0.5cm}}
\dsonly{%
\begin{minipage}[t]{0.3\textwidth}
	\vspace{0pt}\raggedright
	\begin{tcolorbox}[colback=grey]

		\section{Key Skills}
		\begin{itemize}[leftmargin=*]
			\item Machine learning \& time-series forecasting
			\item Statistical modelling
			\item Data engineering \& RAP
			\item CI/CD \& deployment
			\item Graph-based simulation
			\item Data exploration \& visualisation
		\end{itemize}

		\subsection{Programming:}
		\begin{itemize}[leftmargin=*]
			\item Python
			\item SQL
			\item Julia
			\item R
			\item Rust
			\item C/C++
			\item TypeScript
			\item Java
		\end{itemize}

		\subsection{Libraries / Frameworks:}
		\begin{itemize}[leftmargin=*]
			\item numpy / pandas / scikit-learn
			\item PyTorch / TensorFlow
			\item matplotlib / seaborn
			\item React / Next.js
		\end{itemize}

		\subsection{Tools:}
		\begin{itemize}[leftmargin=*]
			\item Git / GitHub
			\item GNU/Linux
			\item Azure DevOps
			\item PostgreSQL / MongoDB
			\item Power BI / Excel
		\end{itemize}

		\subsection{Languages:}
		\begin{itemize}[leftmargin=*]
			\item English (Native)
			\item German (B1)
		\end{itemize}

	\end{tcolorbox}
\end{minipage}
}

\dsonly{%
	\section{Publications}
	\publist
}

\rsonly{%
	\section{Skills}
	\begin{itemize}
		\item \textbf{Research:} tensor networks (MPS), quantum-inspired and probabilistic modelling, time-series forecasting, ML, model evaluation
		\item \textbf{Programming:} Python (numpy, pandas, scikit-learn, PyTorch), Julia (ITensors.jl), SQL, R, Rust, C/C++
		\item \textbf{Tools:} Git, GNU/Linux, PostgreSQL, \LaTeX \quad \textbf{Languages:} English (native), German (B1)
	\end{itemize}
}

\end{document}
